\subsection{The canonical injective resolution}

Let F be the A-module \(\mathrm{Hom}_{\mathbf{Z}}(\mathbf{A},\mathbf{Q} / \mathbf{Z})\) ; for every A-module M, we set \(\mathrm{I}^0 (\mathbf{M}) = \mathrm{F}^{\mathrm{Hom}_{\mathbf{A}}(\mathbf{M},\mathbf{F})}\) and we denote by \(e_{\mathbf{M}}:\mathbf{M}\to \mathrm{I}^{0}(\mathbf{M})\) the homomorphism which maps \(m\in \mathbf{M}\) to the family \((\varphi (m))_{\varphi \in \mathrm{Hom}_{\mathbf{A}}(\mathbf{M},\mathbf{F})}\) . By X, p.~19, cor. 2, \(\mathrm{I}^0 (\mathbf{M})\) is an injective A-module and \(e_{\mathbf{M}}\) is injective. Set \(\mathbf{K}^0 (\mathbf{M}) = \mathbf{Coker}\) \(e_{\mathbf{M}}\) and denote by \(q_{\mathbf{M}}:\mathrm{I}^0 (\mathbf{M})\to \mathrm{K}^0 (\mathbf{M})\) the canonical projection. We thus have an exact sequence

\[
0 \longrightarrow \mathrm{M} \stackrel {e _ {\mathrm{M}}} {\longrightarrow} \mathrm{I} ^ {0} (\mathrm{M}) \stackrel {q _ {\mathrm{M}}} {\longrightarrow} \mathrm{K} ^ {0} (\mathrm{M}) \longrightarrow 0.
\]

We define a graded A-module I(M) by setting \(\mathbf{I}^n (\mathbf{M}) = 0\) for \(n < 0\) and, by induction on the integer \(n > 0\) ,

\[
\mathrm{I} ^ {n} (\mathbf {M}) = \mathrm{I} ^ {0} \left(\mathrm{K} ^ {n - 1} (\mathbf {M})\right), \quad \mathrm{K} ^ {n} (\mathbf {M}) = \mathrm{K} ^ {0} \left(\mathrm{K} ^ {n + 1} (\mathbf {M})\right).\tag{9}
\]

We define A-homomorphisms \(\delta_{\mathbf{M}}^{n}:\mathrm{I}^{n}(\mathbf{M})\to \mathrm{I}^{n + 1}(\mathbf{M})\) by

\[
\left\{ \begin{array}{l l} \delta_ {\mathbf {M}} ^ {n} = 0, & n <   0, \\ \delta_ {\mathbf {M}} ^ {0} = e _ {\mathbf {K} ^ {0} (\mathbf {M})} \circ q _ {\mathbf {M}}, \\ \delta_ {\mathbf {M}} ^ {n} = e _ {\mathbf {K} ^ {n} (\mathbf {M})} \circ q _ {\mathbf {K} ^ {n - 1} (\mathbf {M})}, & n > 0. \end{array} \right.\tag{10}
\]

By construction, we have an exact sequence

\[
0 \longrightarrow \mathrm{M} \xrightarrow {e _ {\mathrm{M}}} \mathrm{I} ^ {0} (\mathrm{M}) \xrightarrow {\delta_ {\mathrm{M}} ^ {0}} \dots \longrightarrow \mathrm{I} ^ {n} (\mathrm{M}) \xrightarrow {\delta_ {\mathrm{M}} ^ {n}} \mathrm{I} ^ {n + 1} (\mathrm{M}) \longrightarrow \dots ,
\]

so that, if one extends \(e_{M}\) as a morphism of complexes

\[
e _ {\mathbf {M}}: \mathbf {M} \rightarrow (\mathrm{I} (\mathbf {M}), \delta_ {\mathbf {M}}),
\]

we obtain an injective resolution of M, called the canonical injective resolution of M. Let \(f: \mathbf{M} \to \mathbf{N}\) be a homomorphism of A-modules. Denote by \(\mathrm{I}^{0}(f)\) the homomorphism from \(\mathrm{I}^{0}(\mathbf{M}) = \mathrm{F}^{\mathrm{Hom}_{\mathbf{A}}(\mathbf{M},\mathbf{F})}\) to \(\mathrm{I}^{0}(\mathbf{N}) = \mathrm{F}^{\mathrm{Hom}_{\mathbf{A}}(\mathbf{N},\mathbf{F})}\) which maps the family \((x_{\varphi})_{\varphi \in \mathrm{Hom}_{\mathbf{A}}(\mathbf{M},\mathbf{F})}\) to the family \((x_{\psi \circ f})_{\psi \in \mathrm{Hom}_{\mathbf{A}}(\mathbf{N},\mathbf{F})}\) . We have:

\[
\mathrm{I} ^ {0} (f) \circ e _ {\mathbf {M}} = e _ {\mathbf {N}} \circ f.\tag{11}
\]

Consequently, \(\mathrm{I}^0 (f)\) induces a homomorphism \(\mathbf{K}^0 (f):\mathbf{K}^0 (\mathbf{M})\to \mathbf{K}_-^0 (\mathbf{N})\) and we have

\[
\mathrm{K} ^ {0} (f) \circ q _ {\mathbf {M}} = q _ {\mathbf {N}} \circ \mathrm{K} ^ {0} (f).\tag{12}
\]

Set \(\mathrm{I}^n(f) = 0\) for \(n < 0\) and let us define, by induction on the integer \(n > 0\), homomorphisms \(\mathrm{I}^n(f): \mathrm{I}^n(\mathbf{M}) \to \mathrm{I}^n(\mathbf{N})\) and \(\mathrm{K}^n(f): \mathrm{K}^n(\mathbf{M}) \to \mathrm{K}^n(\mathbf{N})\) by:

\[
\left\{ \begin{array}{l} \mathrm{I} ^ {n} (f) = \mathrm{I} ^ {0} (\mathrm{K} ^ {n - 1} (f)) \\ \mathrm{K} ^ {n} (f) = \mathrm{K} ^ {0} (\mathrm{K} ^ {n - 1} (f)). \end{array} \right.\tag{13}
\]

PROPOSITION 5. --- I(f) : I(M) → I(N) is a morphism of complexes of A-modules; we have I(f) ∘ e\_M = e\_N ∘ f.

This is proved in a manner analogous to prop. 4.

One has

\[
\mathrm{I} (1 _ {\mathbf {M}}) = 1 _ {\mathrm{I} (\mathbf {M})}\tag{14}
\]

and for every homomorphism \(g: N \rightarrow P\) of A-modules

\[
\mathrm{I} (g \circ f) = \mathrm{I} (g) \circ \mathrm{I} (f).\tag{15}
\]

Remark. --- If \(f, g \in \operatorname{Hom}_{\mathbf{A}}(\mathbf{M}, \mathbf{N})\) , we do not have \(\mathrm{I}(f + g) = \mathrm{I}(f) + \mathrm{I}(g)\) . However, these two morphisms are homotopic by X, p.~49, prop. 3 bis.

If M is a right A-module, we set \(\mathrm{I}(\mathbf{M}) = \mathrm{I}(\mathbf{M}^{\circ})^{\circ}\) ; it is called the canonical injective resolution of M and we have

\[
\mathrm{I} (\mathbf {M} ^ {\circ}) = \mathrm{I} (\mathbf {M}) ^ {\circ}.
\]

